\chapter{Frontend Design}
\label{ch:frontend}

\section{Application shell}

The client is a React 19 single-page application built with Vite, delivered as an
installable progressive web app. Fifteen routes are declared, of which thirteen are
lazy-loaded behind an authentication boundary, so the initial payload contains the landing
page and the shell and nothing else.

\begin{table}[htbp]
\centering
\caption{The route map, the surface each route serves, and its guard.}
\label{tab:routes}
\footnotesize
\begin{tabularx}{\textwidth}{@{}p{3.5cm}p{3.3cm}X@{}}
\toprule
\textbf{Route} & \textbf{Surface} & \textbf{Notes} \\
\midrule
\code{/} & Landing & Public; explains the Circle-First idea before sign-in \\
\code{/enter} & Identity & Nickname, neutral avatar, context; optional account upgrade \\
\code{/dna} & Diagnostic & Four items; writes the passport \\
\code{/solo} & Solo practice & Teach Nova; same ladder, one learner \\
\code{/circles/new}, \code{/circles/join} & Group entry & Create with a window slider; join by code or QR \\
\code{/c/:id} & Lobby & Members, join code, window length, start \\
\code{/c/:id/s/:sid} & Session room & The main surface: window, peers, hints, polls \\
\code{/c/:id/s/:sid/debrief} & Debrief & Final S-B-U, credits, capsule, quiz \\
\code{/board} & Facilitator board & Per-circle rings and \code{needs\_human} \\
\code{/mentor} & Mentor seat & Observer join, note, brief \\
\code{/dashboard} & Analytics & Aggregate-only progress \\
\code{/lab} & Chaos console & Development builds only, excluded from the public bundle \\
\code{/me} & Account & Passport view, export, deletion \\
\bottomrule
\end{tabularx}
\end{table}

Three routing decisions are load-bearing.

\paragraph{Guarded by default.} The API is default-deny, so an anonymous deep link would
render a wall of 401s. Every route past the entry screen is wrapped in a guard that
redirects to \code{/enter} and \emph{remembers the intended destination}, so the flow
resumes after sign-in rather than dumping the learner on the landing page.

\paragraph{The chaos console is a build flag, not a role.} A role check would still ship
the console in the bundle, where it could be reached by anyone willing to read the
JavaScript. The route is gated on \code{import.meta.env.DEV}, so the production bundle
cannot contain it.

\paragraph{Per-route error boundaries.} A lazy chunk that fails to load, or a render error
in one screen, shows a recovery card instead of a blank page. The boundary is keyed by
pathname, so navigating away clears the failure rather than pinning it to the screen.

\section{The session room}

The session room is where the product either works or does not, and it is the screen with
the most state. It renders five things at once, in a fixed vertical order that follows the
learner's attention rather than the data model: the countdown, the question, the peer
answers, the ladder and the verdict.

\begin{figure}[htbp]
\centering
\begin{tikzpicture}[
  font=\sffamily\scriptsize,
  blk/.style={draw=ink0,fill=white,rounded corners=2pt,align=left,inner sep=5pt,
              text width=7.6cm,line width=0.5pt},
  blkf/.style={blk,fill=fillA},
  blkd/.style={blk,fill=fillB},
  tag/.style={font=\sffamily\tiny,color=ink2},
]
\node[blkd] (count) {\textbf{45\,s} \hfill{\color{ink2}peers first}\\[2pt]
  {\color{ink1}Ask: \emph{``Why does binary search need a sorted list?''}}};
\node[blk,below=4mm of count] (peers) {\textbf{Peer answers} \hfill{\color{ink2}2 of 4 answered}\\[2pt]
  \textbullet\ \emph{``you can binary search anything, it's just slower''}
    \hfill{\color{ink2}\tiny flagged}\\
  \textbullet\ \emph{``if it's not sorted you can discard the half that holds it''}};
\node[blkf,below=4mm of peers] (ladder) {\textbf{Hint ladder} \hfill{\color{ink2}L1 used}\\[2pt]
  L1 nudge \textbullet\ L2 analogy \textbullet\ L3 worked start \textbullet\
  {\color{ink2}L4 locked}};
\node[blk,below=4mm of ladder] (nova) {\textbf{Nova} \hfill{\color{ink2}tier C \textbullet\ 5 citations}\\[2pt]
  {\color{ink1}A sorted array is a precondition, not an optimisation. Halving is only
  valid if the discarded half provably cannot contain the target.}\\
  {\color{ink2}\tiny bs-fr-01 · bs-en-01 · bs-ar-01 · bs-darija-01 · bs-misc-01}};
\node[blkf,below=4mm of nova] (poll) {\textbf{Consensus poll} \hfill{\color{ink2}3 votes}\\[2pt]
  Claim: \emph{``binary search works on any array, just slower''} \hfill
  {\color{ink2}needs-fix 3 : 0 correct}};
\draw[decorate,decoration={brace,amplitude=4pt,mirror},ink3]
  ($(count.north west)+(-0.5,0)$) -- ($(count.south west)+(-0.5,0)$)
  node[midway,left=6pt,rotate=90,font=\sffamily\tiny,color=ink2]{before window closes};
\draw[decorate,decoration={brace,amplitude=4pt,mirror},ink3]
  ($(peers.north west)+(-0.5,0)$) -- ($(peers.south west)+(-0.5,0)$)
  node[midway,left=6pt,rotate=90,font=\sffamily\tiny,color=ink2]{circle decides};
\end{tikzpicture}
\caption{The session room, in reading order. The countdown is the hero; peers occupy the
space between asking and knowing; the ladder is offered before the answer; citations sit
under the answer rather than beside it.}
\label{fig:session}
\end{figure}

\subsection{Countdown as a contract}

The countdown is not decoration. It is the visible half of the server-authoritative window:
the client derives the remaining time from \code{windowEndsAt}, renders it, and disables the
peer-answer affordance at zero. Because the deadline lives on the server, a slow clock or a
backgrounded tab cannot extend the window, and the UI cannot accidentally promise time the
server will not grant.

\subsection{Staged progress instead of fake streaming}

Responses arrive whole. The interface therefore shows the pipeline stages
(\emph{retrieving}, \emph{checking similar}, \emph{writing}, \emph{verifying}) as an ordered
stepper with completed steps marked, and never renders a partial answer that the server has
not sent. This is a deliberate honesty choice: a token-by-token animation over a whole
response would be theatre, and the first time it desynchronised from the data it would read
as a bug.

\subsection{The zero-model skeleton}

The first paint of any AI-bearing section is a skeleton that requires no model call. It
renders in under 200\,ms, which is the number in \code{config/router.yaml}, and it is the
reason the screen never shows an empty box while the queue drains. The skeleton is shaped
like the content it will hold, so the transition does not reflow the page.

\section{Design system}

The interface is built from tokens rather than from components with hardcoded values, so
that a change to spacing or type propagates rather than being hunted.

\begin{table}[htbp]
\centering
\caption{Design tokens and the constraint each encodes.}
\label{tab:tokens}
\footnotesize
\begin{tabularx}{\textwidth}{@{}p{3.1cm}X@{}}
\toprule
\textbf{Token group} & \textbf{Rule, and why} \\
\midrule
Type & A single family with a display weight for headings and a text weight for body.
  Optical sizing and tracking are set per size, because a heading tracked for 14\,pt reads
  loose at 32\,pt \\
Space & A 4\,pt base scale. Every margin and padding is a multiple, so vertical rhythm
  holds across screens written by different people \\
Colour & A restrained ink scale plus one accent per context mode. Contrast is asserted
  against WCAG AA in a test, not eyeballed \\
Motion & Durations and easings are tokens; every transition respects
  \code{prefers-reduced-motion} through a single hook \\
Radius and rule & One radius for controls, one for cards, one hairline weight. The
  restraint is what makes dense screens readable \\
Touch & 44\,px minimum target, asserted by a sizing test rather than assumed from padding \\
\bottomrule
\end{tabularx}
\end{table}

\paragraph{No-shame styling.} A design rule with ethical content: a low S-B-U tier is
rendered as a gentle pulse, never as red, and a wrong answer is met with a guiding question
rather than a reprimand. This is implemented at the token level --- there is no ``error red''
in the palette for learner-facing surfaces --- so a future contributor cannot introduce
shame styling without changing the design system.

\paragraph{Anonymous by default where it matters.} Questions can be posted anonymously, and
``me too'' is a private signal between the learner and the coach rather than a public
confession. Both exist because the social cost of asking the simple question is the
specific failure the product is designed to remove.

\section{Internationalization and right-to-left}
\label{sec:i18n}

The interface ships 357 translated keys across four registers: Arabic, Tunisian Darija,
French and English. Three properties were required and are enforced.

\begin{enumerate}[leftmargin=1.6em,itemsep=2pt]
  \item \textbf{Parity.} A checker script walks the key table and fails the build if any key
        is missing a translation in any language, so no screen can fall back to English
        mid-flow.
  \item \textbf{Direction.} Arabic and Darija are right-to-left. Layout uses logical
        properties rather than left/right, so the mirror is a document-level attribute
        change rather than a second stylesheet.
  \item \textbf{Darija is a first-class register, not a translation of Arabic.} The Darija
        strings are written in Darija, including Arabizi-tolerant phrasing where the
        product speaks in that register, and the tests treat Darija as a required locale
        rather than an optional one.
\end{enumerate}

The i18n test also asserts that no user-visible string is hardcoded in a component ---
a rule that is easy to state and easy to break under deadline, which is exactly why it is
a test.

\section{Accessibility}

Accessibility is treated as a set of gates rather than a checklist. The gates that run on
every build are:

\begin{itemize}[leftmargin=1.4em,itemsep=2pt]
  \item \textbf{axe-core on rendered screens}, including the 404 and the crash card. The
        gate runs the real engine against the real DOM rather than asserting hand-counted
        attributes, so a component that loses its label fails.
  \item \textbf{Keyboard completeness}: every interactive element is reachable and
        operable, and focus is visible. A focus trap in a sheet is tested rather than
        trusted.
  \item \textbf{Announced route transitions}: the lazy-chunk fallback carries a status role
        and a label in the learner's own language, so a screen reader announces
        ``loading'' in Arabic rather than in English.
  \item \textbf{Reduced motion}: one hook owns the media query, so no component can opt out
        of it by accident.
\end{itemize}

\begin{note}[Why the crash card is tested for accessibility]
The screen a user sees when something breaks is the screen they see at their most
confused. It is also the screen most likely to be written in a hurry. Testing it for
contrast, focus and labels catches the class of bug that only appears in the worst
moment.
\end{note}

\section{Offline behaviour}

The service worker precaches the application shell and the offline pack, so a dropped
connection does not present a browser error page. In the \code{BLACK} degradation state the
client serves the pack locally with an explicit \code{offline:true} label, which is the
honest version of the promise: the learner knows the answer came from a snapshot rather
than from the model. Writes that arrive while offline are queued and synced later, and the
session continues with peer work, polls, role-play and flashcards --- the activities that
never needed the model in the first place.

\section{Performance}

The performance budget is small and stated: a skeleton in under 200\,ms, first paint from
the precached shell on a repeat visit, and no model call on the critical path of a route
transition. Route-level code splitting keeps the initial bundle to the landing page plus the
shell; the session room, which is the heaviest screen, is fetched only when a learner enters
a circle.
